\BeleskeID{02-s015b}
% element: 02-s015b:e01
% element: 02-s015b:e02
% element: 02-s015b:e03
% element: 02-s015b:d01

Za preostale redove 010, 100 i 111 dobijamo redom \(C+\bar B+A\), \(\bar C+B+A\) i \(\bar C+\bar B+\bar A\). Svaki od tih potpunih zbirova jednak je nuli samo u odgovarajućem redu. Međukorak koji povezuje svih pet uslova glasi
\[
\begin{aligned}
F={}&\overline{\bar C\bar B\bar A}\cdot\overline{\bar C\bar BA}\cdot\overline{\bar CB\bar A}\\
 &\cdot\overline{C\bar B\bar A}\cdot\overline{CBA}.
\end{aligned}
\]
Primenom De Morganovog pravila na svaki činilac dobijamo proizvod potpunih zbirova. Ne sabiramo te zbirove: dovoljno je da jedan od njih bude 0 kako bi konačan izlaz bio 0. U redovima 011, 101 i 110 nijedan uslov za nulu nije ispunjen, pa svih pet činilaca iznosi 1 i proizvod daje 1.
